\tagstructbegin{tag=Figure,alttext={两条计算路径。教学输入 ID 2、4、5 进入五项词表的三维教学模型，得到三行五列概率。报告原文与任务指令进入能处理报告的预训练语言模型，得到结构化提取草稿，再由程序核对出处与环比。两种模型的计算原理相通，使用各自的参数；处理报告请求时真实模型参数固定。}}
\begin{center}
\begin{tikzpicture}[>=Latex,x=1mm,y=1mm,
 n/.style={draw=linegray,rounded corners=2pt,fill=paper,align=center,
 font=\sffamily\small,minimum height=17mm,text width=29mm}]
\node[n] (toyinput) at (0,0) {教学输入\\ID $[2,4,5]$};
\node[n,fill=linegray!15] (toy) at (46,0) {三维教学模型\\五项词表};
\node[n] (probs) at (92,0) {三个位置的概率\\$\Pi:3\times5$};
\draw[->,thick,flowcolor] (toyinput)--(toy);
\draw[->,thick,flowcolor] (toy)--(probs);
\node[n] (report) at (0,-37) {报告原文＋任务\\含原文出处};
\node[n,fill=linegray!15] (real) at (46,-37) {预训练语言模型\\请求中 $\theta_{\rm real}$ 固定};
\node[n] (draft) at (92,-37) {结构化提取草稿\\季度、收入、出处};
\node[n] (check) at (92,-64) {程序核对出处\\计算 $20\%$、$10\%$};
\draw[->,thick,flowcolor] (report)--(real);
\draw[->,thick,flowcolor] (real)--(draft);
\draw[->,thick,flowcolor] (draft)--(check);
\draw[densely dotted,thick,draw=muted] (toy)--node[right,font=\footnotesize,text=muted]{计算原理相通} (real);
\end{tikzpicture}
\captionof{figure}{从教学算例的计算原理到真实模型的报告应用}
\label{fig:model-to-assistant}
\end{center}
\tagstructend
